% GENERATED FILE. Edit canonical lessons, then run npm run generate:books.
% canonical-source-hash: fnv1a64:eef2a1653e1f897b
% canonical-lessons: IT-C475-spavento, IT-C475-spaventarsi, IT-C475-panico, IT-C475-sollievo, IT-C475-tranquillizzarsi

\chapter{Che spavento!}
\label{ch:it-che-spavento}
\hlchaptermodality{475}

\begin{chapteropening}
\textbf{By the end of this chapter:} \emph{I can talk about a fright and the relief afterwards.}

Complete the last lesson of chapter 475: I can talk about a fright and the relief afterwards.
\end{chapteropening}

\section[lo spavento]{lo spavento — fright, scare}

\label{lesson:IT-C475-spavento}



\begin{quote}
Before the new one: say the Italian for to get irritated, then the Italian for to huff.
\end{quote}

\subsection*{You'll want to know: lo spavento}
\textbf{lo spavento} — \textquotedblleft{}fright, scare\textquotedblright{}.

\begin{quote}
Che spavento! Credevo di aver perso le chiavi.
\end{quote}

\begin{quote}
\emph{What a fright! I thought I had lost the keys.}
\end{quote}

\emph{Prendersi uno spavento} — to get a fright.

\begin{cousinweb}
\emph{spavento} comes from \emph{spaventare} \textquotedblleft{}to frighten\textquotedblright{}, usually traced to Vulgar Latin \emph{expaventāre}, built on Latin \emph{pavēre} \textquotedblleft{}to be afraid\textquotedblright{} (Lewis \& Short).
\end{cousinweb}

\begin{grammarlens}[title={What you've built}]
One more word for telling, explaining and describing at B1.
\end{grammarlens}

\subsection*{Your turn}
\begin{itemize}
\raggedright
  \item \emph{Say it:} \emph{lo spavento}
  \item \emph{Say it:} \emph{lo spavento}, once more
  \item \emph{Say it:} a sentence of your own with \emph{lo spavento}
  \item \emph{Recall:} say the Italian for to get irritated, then the Italian for to huff, then say \emph{lo spavento} again
\end{itemize}

\begin{tcolorbox}[breakable,colback=teal!4,colframe=teal!35!black,title={Before you move on}]
What does lo spavento mean? (\textquotedblleft{}Fright, scare\textquotedblright{}.) Say it once more.
\end{tcolorbox}
\section[spaventarsi]{spaventarsi — to get frightened}

\label{lesson:IT-C475-spaventarsi}



\begin{quote}
Before the new one: say the Italian for to huff, then the Italian for fright.
\end{quote}

\subsection*{You'll want to know: spaventarsi}
\textbf{spaventarsi} — \textquotedblleft{}to get frightened\textquotedblright{}.

\begin{quote}
Il cane si spaventa quando sente i tuoni.
\end{quote}

\begin{quote}
\emph{The dog gets frightened when it hears thunder.}
\end{quote}

Without \emph{si}, \emph{spaventare} is to frighten someone: \emph{mi hai spaventato!}

\begin{cousinweb}
\emph{spaventarsi} is the reflexive of \emph{spaventare} \textquotedblleft{}to frighten\textquotedblright{}, usually traced to Vulgar Latin \emph{expaventāre}, from Latin \emph{pavēre} \textquotedblleft{}to be afraid\textquotedblright{} (Lewis \& Short).
\end{cousinweb}

\begin{grammarlens}[title={What you've built}]
One more word for telling, explaining and describing at B1.
\end{grammarlens}

\subsection*{Your turn}
\begin{itemize}
\raggedright
  \item \emph{Say it:} \emph{spaventarsi}
  \item \emph{Say it:} \emph{spaventarsi}, once more
  \item \emph{Say it:} a sentence of your own with \emph{spaventarsi}
  \item \emph{Recall:} say the Italian for to huff, then the Italian for fright, then say \emph{spaventarsi} again
\end{itemize}

\begin{tcolorbox}[breakable,colback=teal!4,colframe=teal!35!black,title={Before you move on}]
What does spaventarsi mean? (\textquotedblleft{}To get frightened\textquotedblright{}.) Say it once more.
\end{tcolorbox}
\section[il panico]{il panico — panic}

\label{lesson:IT-C475-panico}



\begin{quote}
Before the new one: say the Italian for fright, then the Italian for to get frightened.
\end{quote}

\subsection*{You'll want to know: il panico}
\textbf{il panico} — \textquotedblleft{}panic\textquotedblright{}.

\begin{quote}
Non farti prendere dal panico, troviamo una soluzione.
\end{quote}

\begin{quote}
\emph{Do not panic, we will find a solution.}
\end{quote}

\emph{Farsi prendere dal panico} — to panic.

\begin{cousinweb}
\emph{panico} comes through Latin \emph{pānicus} from Greek \emph{panikós} (EtymDB), an adjective made from the name of a woodland god whose voice was said to fill travellers with sudden fear.
\end{cousinweb}

\begin{grammarlens}[title={What you've built}]
One more word for telling, explaining and describing at B1.
\end{grammarlens}

\subsection*{Your turn}
\begin{itemize}
\raggedright
  \item \emph{Say it:} \emph{il panico}
  \item \emph{Say it:} \emph{il panico}, once more
  \item \emph{Say it:} a sentence of your own with \emph{il panico}
  \item \emph{Recall:} say the Italian for fright, then the Italian for to get frightened, then say \emph{il panico} again
\end{itemize}

\begin{tcolorbox}[breakable,colback=teal!4,colframe=teal!35!black,title={Before you move on}]
What does il panico mean? (\textquotedblleft{}Panic\textquotedblright{}.) Say it once more.
\end{tcolorbox}
\section[il sollievo]{il sollievo — relief}

\label{lesson:IT-C475-sollievo}



\begin{quote}
Before the new one: say the Italian for to get frightened, then the Italian for panic.
\end{quote}

\subsection*{You'll want to know: il sollievo}
\textbf{il sollievo} — \textquotedblleft{}relief\textquotedblright{}.

\begin{quote}
Che sollievo, sei arrivato!
\end{quote}

\begin{quote}
\emph{What a relief, you have arrived!}
\end{quote}

\emph{Tirare un sospiro di sollievo} — to breathe a sigh of relief.

\begin{cousinweb}
\emph{sollievo} is built on \emph{sollevare} \textquotedblleft{}to lift up\textquotedblright{}, which is inherited from Latin \emph{sublevāre} \textquotedblleft{}to lift from below\textquotedblright{} (Lewis \& Short), from \emph{levāre} \textquotedblleft{}to raise\textquotedblright{}.
\end{cousinweb}

\begin{grammarlens}[title={What you've built}]
One more word for telling, explaining and describing at B1.
\end{grammarlens}

\subsection*{Your turn}
\begin{itemize}
\raggedright
  \item \emph{Say it:} \emph{il sollievo}
  \item \emph{Say it:} \emph{il sollievo}, once more
  \item \emph{Say it:} a sentence of your own with \emph{il sollievo}
  \item \emph{Recall:} say the Italian for to get frightened, then the Italian for panic, then say \emph{il sollievo} again
\end{itemize}

\begin{tcolorbox}[breakable,colback=teal!4,colframe=teal!35!black,title={Before you move on}]
What does il sollievo mean? (\textquotedblleft{}Relief\textquotedblright{}.) Say it once more.
\end{tcolorbox}
\section[tranquillizzarsi]{tranquillizzarsi — to calm down}

\label{lesson:IT-C475-tranquillizzarsi}



\begin{quote}
Before the new one: say the Italian for panic, then the Italian for relief.
\end{quote}

\subsection*{You'll want to know: tranquillizzarsi}
\textbf{tranquillizzarsi} — \textquotedblleft{}to calm down\textquotedblright{}.

\begin{quote}
Respira piano e tranquillizzati.
\end{quote}

\begin{quote}
\emph{Breathe slowly and calm down.}
\end{quote}

Built on \emph{tranquillo}. To calm someone else: \emph{tranquillizzare qualcuno}.

\begin{cousinweb}
\emph{tranquillizzarsi} is built on \emph{tranquillo} \textquotedblleft{}calm\textquotedblright{}, a learned borrowing from Latin \emph{tranquillus} \textquotedblleft{}calm, still\textquotedblright{} (Lewis \& Short), with the verb ending \emph{-izzare}.
\end{cousinweb}

\begin{grammarlens}[title={What you've built}]
One more word for telling, explaining and describing at B1.
\end{grammarlens}

\subsection*{Your turn}
\begin{itemize}
\raggedright
  \item \emph{Say it:} \emph{tranquillizzarsi}
  \item \emph{Say it:} \emph{tranquillizzarsi}, once more
  \item \emph{Say it:} a sentence of your own with \emph{tranquillizzarsi}
  \item \emph{Recall:} say the Italian for panic, then the Italian for relief, then say \emph{tranquillizzarsi} again
\end{itemize}

\begin{tcolorbox}[breakable,colback=teal!4,colframe=teal!35!black,title={Before you move on}]
What does tranquillizzarsi mean? (\textquotedblleft{}To calm down\textquotedblright{}.) Say it once more.
\end{tcolorbox}
